\documentclass[border=5pt]{standalone}
\usepackage[T1]{fontenc}
\usepackage{tikz}
\usepackage{luacode}

\definecolor{garnet}{HTML}{73000A}
\definecolor{atlantic}{HTML}{466A9F}
\definecolor{congaree}{HTML}{1F414D}
\definecolor{warmgrey}{HTML}{676156}
\definecolor{bk70}{HTML}{5C5C5C}

\begin{document}
\begin{tikzpicture}

\begin{luacode}
-- ════════════════════════════════════════════════════════════
-- Animated lens morph: \PARAM sweeps 0..360
-- 6 keyframes at 0, 60, 120, 180, 240, 300
-- Each segment: 20 frames hold + 40 frames morph to next
-- ════════════════════════════════════════════════════════════

local frame = \PARAM

-- Keyframes: {Rl, Rr, label}
-- Flat keyframes inserted between sign changes to avoid R crossing 0
local keys = {
    { 3,   -3,   "Biconvex (symmetric)"},
    { 2,   -6,   "Biconvex (asymmetric)"},
    { 20,  -20,  "Flattening..."},          -- bridge: Rl sign change next
    {-3,    3,   "Biconcave (symmetric)"},
    {-20,   20,  "Flattening..."},           -- bridge: Rl sign change next
    { 3,    3,   "Meniscus (converging)"},
    { 20,  -3,   "Plano-convex"},
    { 20,   3,   "Plano-concave"},
}

local n_keys = #keys
local seg_len = 60       -- frames per segment
local hold = 20          -- frames to hold before morphing
local morph = seg_len - hold  -- frames for transition

-- Figure out which segment and local frame
local seg = math.floor(frame / seg_len)         -- 0..5
local local_f = frame - seg * seg_len            -- 0..59
if seg >= n_keys then seg = n_keys - 1; local_f = seg_len - 1 end

local k1 = keys[seg + 1]
local k2 = keys[math.min(seg + 2, n_keys)]  -- next keyframe (clamp at end)

local Rl, Rr, label
if local_f < hold then
    -- Hold current keyframe
    Rl = k1[1]
    Rr = k1[2]
    label = k1[3]
else
    -- Morph to next
    local t = (local_f - hold) / morph
    -- Smooth step for nicer easing
    t = t * t * (3 - 2 * t)
    Rl = k1[1] + (k2[1] - k1[1]) * t
    Rr = k1[2] + (k2[2] - k1[2]) * t
    -- Show target label once past halfway
    label = (t < 0.5) and k1[3] or k2[3]
end

-- Safety clamp: |R| must stay above h + margin
local min_R = 1.8
if math.abs(Rl) < min_R then Rl = (Rl >= 0 and 1 or -1) * min_R end
if math.abs(Rr) < min_R then Rr = (Rr >= 0 and 1 or -1) * min_R end

-- ════════════════════════════════════════════════════════════
-- Fixed config
-- ════════════════════════════════════════════════════════════
local h  = 1.2
local n_glass = 1.52
local n_rays     = 30
local ray_y_min  = -1.1
local ray_y_max  =  1.1
local ray_width  = 0.2
local ray_opacity = 0.5
local lens_x = 2.5
local xmin   = -0.5
local xmax   = 6.5
local min_wall = 0.08

-- ════════════════════════════════════════════════════════════
-- ENGINE
-- ════════════════════════════════════════════════════════════

local function auto_t(Rl, Rr, h)
    local aRl, aRr = math.abs(Rl), math.abs(Rr)
    local sl = (aRl >= h) and (aRl - math.sqrt(aRl*aRl - h*h)) or aRl
    local sr = (aRr >= h) and (aRr - math.sqrt(aRr*aRr - h*h)) or aRr
    local dl = (Rl > 0) and -sl or sl
    local dr = (Rr < 0) and -sr or sr
    local t = min_wall - math.min(0, dl + dr)
    if t < min_wall then t = min_wall end
    return t
end

local function isect(ox, oy, dx, dy, cx, cy, r, pick)
    local ex, ey = ox - cx, oy - cy
    local a = dx*dx + dy*dy
    local b = 2*(ex*dx + ey*dy)
    local c = ex*ex + ey*ey - r*r
    local disc = b*b - 4*a*c
    if disc < 0 then return nil end
    local sq = math.sqrt(disc)
    local t1 = (-b - sq)/(2*a)
    local t2 = (-b + sq)/(2*a)
    local t
    if pick == "far" then
        t = (t2 > 1e-4) and t2 or ((t1 > 1e-4) and t1 or nil)
    else
        t = (t1 > 1e-4) and t1 or ((t2 > 1e-4) and t2 or nil)
    end
    if not t then return nil end
    local hx, hy = ox + t*dx, oy + t*dy
    return t, hx, hy, (hx-cx)/r, (hy-cy)/r
end

local function refract(dx, dy, nx, ny, n1, n2)
    local d = dx*nx + dy*ny
    if d > 0 then nx, ny, d = -nx, -ny, -d end
    local r = n1/n2
    local ci = -d
    local s2t = r*r*(1 - ci*ci)
    if s2t > 1 then return nil end
    local ct = math.sqrt(1 - s2t)
    local rx = r*dx + (r*ci - ct)*nx
    local ry = r*dy + (r*ci - ct)*ny
    local l = math.sqrt(rx*rx + ry*ry)
    return rx/l, ry/l
end

local function lens_geo(lx, Rl, Rr, h)
    local t = auto_t(Rl, Rr, h)
    local vl = lx - t/2
    local vr = lx + t/2
    return { vl=vl, vr=vr, t=t,
             cl_x = vl + Rl, cr_x = vr + Rr,
             aRl=math.abs(Rl), aRr=math.abs(Rr) }
end

local function draw_lens(lx, Rl, Rr, h)
    local G = lens_geo(lx, Rl, Rr, h)
    local N = 80
    local lp, rp = {}, {}
    for i = 0, N do
        local y = -h + 2*h*(i/N)
        local ql = math.sqrt(math.max(G.aRl*G.aRl - y*y, 0))
        local qr = math.sqrt(math.max(G.aRr*G.aRr - y*y, 0))
        table.insert(lp, {(Rl > 0) and (G.cl_x - ql) or (G.cl_x + ql), y})
        table.insert(rp, {(Rr < 0) and (G.cr_x + qr) or (G.cr_x - qr), y})
    end
    local p = "\\fill[atlantic!12, opacity=0.85] "
    for i, v in ipairs(lp) do
        p = p .. (i==1 and "" or " -- ") .. string.format("(%.4f,%.4f)", v[1], v[2])
    end
    for i = #rp, 1, -1 do
        p = p .. string.format(" -- (%.4f,%.4f)", rp[i][1], rp[i][2])
    end
    tex.sprint(p .. " -- cycle;")
    for _, arc in ipairs({lp, rp}) do
        local s = "\\draw[atlantic!60, thin] "
        for i, v in ipairs(arc) do
            s = s .. (i==1 and "" or " -- ") .. string.format("(%.4f,%.4f)", v[1], v[2])
        end
        tex.sprint(s .. ";")
    end
    tex.sprint(string.format("\\draw[atlantic!60, thin] (%.4f,%.4f)--(%.4f,%.4f);",
        lp[1][1], lp[1][2], rp[1][1], rp[1][2]))
    tex.sprint(string.format("\\draw[atlantic!60, thin] (%.4f,%.4f)--(%.4f,%.4f);",
        lp[#lp][1], lp[#lp][2], rp[#rp][1], rp[#rp][2]))
end

-- Clip ray endpoint to bounding box, return the point where it exits
local bbox = {xmin = xmin, xmax = xmax, ymin = -1.8, ymax = 2.5}
local function clip_to_bbox(cx, cy, cdx, cdy)
    local t_best = 1e9
    -- Check all 4 walls
    if math.abs(cdx) > 1e-9 then
        local t = (bbox.xmax - cx) / cdx
        if t > 0 and t < t_best then t_best = t end
        t = (bbox.xmin - cx) / cdx
        if t > 0 and t < t_best then t_best = t end
    end
    if math.abs(cdy) > 1e-9 then
        local t = (bbox.ymax - cy) / cdy
        if t > 0 and t < t_best then t_best = t end
        t = (bbox.ymin - cy) / cdy
        if t > 0 and t < t_best then t_best = t end
    end
    if t_best < 1e9 then
        return cx + t_best * cdx, cy + t_best * cdy
    end
    return cx, cy
end

local function trace(ox, oy, dx, dy, lx, Rl, Rr, h)
    local G = lens_geo(lx, Rl, Rr, h)
    local pts = {{ox, oy}}
    local cx, cy, cdx, cdy = ox, oy, dx, dy
    local pick_l = (Rl > 0) and "near" or "far"
    local t, hx, hy, nx, ny = isect(cx, cy, cdx, cdy, G.cl_x, 0, G.aRl, pick_l)
    if t and math.abs(hy) <= h then
        table.insert(pts, {hx, hy})
        local ndx, ndy = refract(cdx, cdy, nx, ny, 1.0, n_glass)
        if not ndx then
            local ex, ey = clip_to_bbox(hx, hy, cdx, cdy)
            table.insert(pts, {ex, ey}); return pts
        end
        cdx, cdy = ndx, ndy
        cx, cy = hx, hy
        local pick_r = (Rr < 0) and "far" or "near"
        t, hx, hy, nx, ny = isect(cx, cy, cdx, cdy, G.cr_x, 0, G.aRr, pick_r)
        if t and math.abs(hy) <= h then
            table.insert(pts, {hx, hy})
            ndx, ndy = refract(cdx, cdy, nx, ny, n_glass, 1.0)
            if not ndx then
                local ex, ey = clip_to_bbox(hx, hy, cdx, cdy)
                table.insert(pts, {ex, ey}); return pts
            end
            cdx, cdy = ndx, ndy
            cx, cy = hx, hy
        end
    end
    -- Extend to bounding box edge
    local ex, ey = clip_to_bbox(cx, cy, cdx, cdy)
    table.insert(pts, {ex, ey})
    return pts
end

-- ══════════ DRAW ══════════

-- Optical axis
tex.sprint(string.format(
    "\\draw[warmgrey, dashed, thin] (%.2f,0)--(%.2f,0);", xmin, xmax))

-- Lens
draw_lens(lens_x, Rl, Rr, h)

-- Rays
for i = 0, n_rays - 1 do
    local frac = (n_rays > 1) and (i / (n_rays - 1)) or 0.5
    local y0 = ray_y_min + (ray_y_max - ray_y_min) * frac
    local pts = trace(xmin, y0, 1, 0, lens_x, Rl, Rr, h)
    for j = 1, #pts - 1 do
        tex.sprint(string.format(
            "\\draw[garnet, line width=%.2fpt, opacity=%.2f] (%.4f,%.4f)--(%.4f,%.4f);",
            ray_width, ray_opacity,
            pts[j][1], pts[j][2], pts[j+1][1], pts[j+1][2]))
    end
end

-- Labels
local rl_disp = string.format("%.1f", Rl)
local rr_disp = string.format("%.1f", Rr)
tex.sprint(string.format(
    "\\node[font=\\sffamily\\small, text=bk70, anchor=south] at (%.2f,%.2f) {$R_L = %s$ \\qquad $R_R = %s$};",
    lens_x, h + 0.35, rl_disp, rr_disp))

-- Escape label for TeX
local safe_label = label:gsub("%(", "{\\lparen}"):gsub("%)", "{\\rparen}")
tex.sprint(string.format(
    "\\node[font=\\sffamily\\footnotesize\\bfseries, text=atlantic, anchor=south] at (%.2f,%.2f) {%s};",
    lens_x, h + 0.7, label))

\end{luacode}

\useasboundingbox (-0.5,-1.8) rectangle (6.5,2.5);
\clip (-0.5,-1.8) rectangle (6.5,2.5);
\draw[bk70, very thin] (-0.5,-1.8) rectangle (6.5,2.5);

\end{tikzpicture}
\end{document}
